\todo[inline, color=pink, size=normalsize]{Descrição da abordagem utilizada (estudos de caso foram selecionados e resolvidos com vários métodos e pacotes)}

O presente capítulo tem como objetivo apresentar a metodologia proposta para a resolução de PCOs, bem como destacar os estudos de caso que serão avaliados com a mesma. 

\section{\boldmath$ COPILOTS $}

O $COPILOTS$ (\textit{Basic OptimaL Control Solver}) é um pacote desenvolvido neste trabalho para resolução de PCOs a partir da implementação de Métodos de Colocação Direta, mais especificamente da colocação trapezoidal e da colocação Hermite-Simpson. Esse pacote, escrito para o Maltab\textsuperscript{\textregistered}, foi elaborado para ser utilizado por usuários com pouca ou nenhuma experiência na resolução computacional de PCOs. Para isso, o $COPILOTS$ possibilita que, por meio da execução de um único comando, sejam criados, e já parcialmente preenchidos, os \textit{scripts} a serem utilizados na estruturação do PCO, guiando o usuário iniciante. Além disso, o $COPILOTS$ apresenta uma sintaxe algébrica próxima à da formulação de Bolza, o que simplifica a implementação do PCO \cite{febbo_nloptcontrol_2020}. O PPNL resultante da discretização é resolvido via uso do algoritmo SQP (\textit{Sequential Quadratic Programming}) \cite{vanderplaats_numerical_1984}.	

Em termos comparativos, considerando um das métricas propostas em \citeonline{febbo_nloptcontrol_2020} (o número de linhas de código), a implementação do estudo de caso apresentado na Seção \ref{sec:resultados:estacionamento}, por exemplo, é mais simples no $ COPILOTS $ do que em outros pacotes similares, conforme observado na Tab. \ref{tab:revisao:comparacaoLinhas}. Vale ressaltar que o $ FALCON $ e o $ ICLOCS $ (\textit{Imperial College London Optimal Control Software}) \cite{falugi_iclocs2_2018} são pacotes escritos para o Matlab\textsuperscript{\textregistered}, enquanto o $ PSOPT $ é baseado em C++. 

\begin{table}[!htb]
	\caption{Número de linhas de código utilizadas na implementação do estudo de caso introduzido na Seção \ref{sec:resultados:estacionamento}.}
	\label{tab:revisao:comparacaoLinhas}
	\centering{}
	\begin{tabular}{|c|c|}
		\hline
		$COPILOTS$  & 149 \\ \hline
		$ FALCON $ & 165 \\ \hline
		$ PSOPT $    & 311 \\ \hline
		$ ICLOCS $   & 362 \\ \hline
	\end{tabular}
\end{table}

No $ COPILOTS $ é possível visualizar e alterar o código para cada estudo de caso. Além disso, o pacote encontra-se estruturado em módulos e foi desenvolvido com base no paradigma da Programação Orientada a Objetos, o que facilita a compreensão do código fonte e a implementação de atualizações \cite{parejo_metaheuristic_2012}.

O processo de instalação do $ COPILOTS $ é composto por duas etapas. Primeiramente, faz-se o \textit{download} do código fonte em \url{https://cutt.ly/wjYNIik} e adiciona-se a pasta em que o mesmo tenha sido armazenado ao \textit{path} do Matlab\textsuperscript{\textregistered}  \cite{mathworks_change_2020}. Em seguida, executa-se no Matlab\textsuperscript{\textregistered} o comando \texttt{copilotsSetup}. 

A resolução de um PCO no $ COPILOTS $ também é um processo constituído de duas etapas. Primeiramente executa-se o comando \texttt{copilotsNew} para criação da pasta \textit{new}, que contém os \textit{scripts} a serem preenchidos para implementação do PCO. Uma vez realizado esse preenchimento, executa-se o comando \texttt{copilots} dentro da pasta \textit{new} para que se inicie o processo de obtenção da solução. Ao fim da execução a pasta \textit{results} será criada para armazenamento dos resultados obtidos, como as trajetórias ótimas dos estados e controles, e os parâmetros associados à resolução do PPNL. A interpolação dos estados e controles é realizada de forma automática e a representação gráfica das trajetórias ótimas é, por padrão, apresentada ao fim da execução.

Por fim vale ressaltar que tanto o código fonte do $COPILOTS$ quanto os arquivos gerados após a execução do comando \texttt{copilotsNew} estão repletos de comentários que orientam a utilização do pacote. Mais ainda, acompanham o código fonte uma série de exemplos que podem ser utilizadas como base para implementação de novos estudos de caso.

\textcolor{red}{Arthur, não seria interessante colocar uma tela de apresentação do pacote, ou algo que o represente???? Um fluxograma ou algo do tipo ....}

\section{Estudos de caso}

A escolha dos estudos de caso para fins de aplicação não é uma tarefa trivial. Em linhas gerais pode-se dizer que o conjunto de aplicações deve ser extenso, heterogêneo \cite{bongartz_numerical_1997} e composto por problemas complexos e interessantes \cite{dolan_benchmarking_2002}. De fato, a complexidade de um PCO pode ser inferida com base no número de estados, controles, e restrições \cite{dolan_benchmarking_2002}, ou ainda a não linearidade da função objetivo e das restrições \cite{bongartz_numerical_1997}. Porém, apenas com base em uma análise subjetiva é possível dizer se um problema é interessante ou não. Assim, não há um consenso a respeito de quais estudos de caso devem ser empregados na elaboração de um estudo comparativo \cite{dolan_benchmarking_2002}. 

Nesse contexto, os estudos de caso abordados no presente trabalho foram escolhidos segundo os critérios de complexidade destacados anteriormente. Estes são descritos como:
%
\begin{enumerate}
\item Minimização do esforço durante a aceleração de um bloco \cite{becerra_optimal_2008};
\item Problemas singulares: Casos 1 e Caso 2 \cite{jacobson_computation_1970};
\item Minimização do esforço durante o \textit{Swing-up} de um pêndulo invertido \cite{kelly_introduction_2017};
\item Minimização do tempo durante uma manobra de estacionamento \cite{li_time-optimal_2016};
\item Otimização da trajetória de um UAV (\textit{Unmanned Aerial Vehicle}) \cite{toledo_de_azevedo_pseudospectral_2018};
\item Lançamento do foguete Delta III \cite{benson_gauss_2005}.
\end{enumerate}

\todo[inline, color=pink, size=normalsize]{Porque escolhi cada estudo de caso do jeito que escolhi}

O estudo de caso (ou problema) 1 possui uma dinâmica bastante simples e poucas restrições, sendo o ponto de partida ideal para um usuário iniciante. Já os problemas singulares são de implementação simples, porém de solução complexa, dadas as propriedades numéricas atribuídas aos mesmos. O problema 3 têm os estados iniciais e finais previamente estabelecidos, enquanto o problema 4 possui um número consideravelmente elevado de restrições e é bastante sensível aos palpites inicias atribuídos aos estados e controles. A dinâmica do problema 5 é descontínua, uma vez que depende da interpolação linear dos dados de uma tabela, e o problema 6 consiste em um PCO de múltiplas fases. 

Vale ressaltar que, enquanto os índices de desempenho associados aos estudos de caso 1, 3 e 5 dependem apenas do termo de Lagrange, que contabiliza a evolução temporal dos estados e/ou controles, os índices de desempenho associados aos estudos de caso 2, 4 e 6 dependem somente do termo de Mayer, computado com base nos valores finais dos estados e/ou do tempo gasto na execução da trajetória. Assim sendo, conclui-se que o conjunto dos estudos de caso selecionados é consideravelmente heterogêneo. 

\textcolor{red}{Arthur, não acho interessante colocarmos apelidos aos problemas. Pore exemplo, podemos chamar o primeiro de Problema da Aceleração de um Bloco, Problema Singular 1, Problema Singular 2, etc .. Por isso, retirei a tabela ...}

%\todo[inline, color=pink, size=normalsize]{Apelidos dos estudos de caso}

%Para que os estudos de caso possam ser referenciados de forma conveniente nas legendas e títulos dos gráficos apresentados no Capítulo \ref{sec:resultados}, atribui-se a cada um desses estudos de caso um apelido. A relação entre os estudos de caso, os apelidos associados aos mesmos, e as seções nas quais são introduzidos, é apresentada na Tabela \ref{tab:metodologia:apelidos}.

%\begin{table}[!htb]
%	\centering
%	\caption[Relação entre os estudos de caso, os apelidos associados aos mesmos, e as seções nas quais são introduzidos]{Relação entre os estudos de caso, os apelidos associados aos mesmos, e as seções nas quais são introduzidos.}
%	\label{tab:metodologia:apelidos}
%	\begin{tabu} to \textwidth {| X[3,m] | X[m,c] | X[m,c] |}
%		\hline
%		\textbf{Estudo de caso} & \vspace{5pt} \textbf{Número da seção} \vspace{5pt} & \textbf{Apelido} \\ \hline
%		\vspace{5pt} Aceleração de um bloco empregando o mínimo esforço \vspace{5pt} & \ref{sec:resultados:integrador} & Bloco \\ \hline
%		\vspace{5pt} Problemas singulares: Caso 1 \vspace{5pt} & \ref{sec:resultados:singular1} & Singular 1 \\ \hline
%		\vspace{5pt} Problemas singulares: Caso 2 \vspace{5pt} & \ref{sec:resultados:singular2} & Singular 2 \\ \hline
%		\vspace{5pt} \textit{Swing-up} de um pêndulo invertido empregando o mínimo esforço \vspace{5pt} & \ref{sec:resultados:penduloInvertido} & Pêndulo Invertido \\ \hline
%		\vspace{5pt} Realização de uma manobra de estacionamento paralelo em tempo mínimo \vspace{5pt} & \ref{sec:resultados:estacionamento} & Estacionamento \\ \hline
%		\vspace{5pt} Otimização da trajetória de um UAV, na presença de um campo de vento, para minimização do consumo de bateria \vspace{5pt} & \ref{sec:resultados:uav} & UAV \\ \hline
%		\vspace{5pt} Lançamento de um foguete Delta III \vspace{5pt} & \ref{sec:resultados:foguete} & Foguete \\ \hline
%	\end{tabu}
%end{table}


\todo[inline, color=pink, size=normalsize]{Métricas utilizadas}

Conforme comentado no capítulo anterior, várias métricas de comparação podem ser empregadas para avaliar pacotes computacionais. No presente neste trabalho serão empregadas as seguintes:
%
\begin{enumerate}
	\item \textbf{Valor ótimo da função objetivo - \boldmath$J^* $}
	
	Consiste na comparação entre o melhor valor obtido para a função objetivo $ J $.
	
	\item \textbf{Tempo de processamento - \boldmath$ t_p $}
	
	Equivale ao tempo despendido na solução do PCO. Uma vez que o sistema operacional empregado na geração dos resultados apresentados no presente trabalho não é um operacional de tempo real, é necessário que se tenha em mente que a cada execução  atribui-se um $ t_p $ distinto. Desta forma,
	para que o $ t_p $ associado à resolução de um dado PCO seja determinado é preciso resolvê-lo várias vezes de forma que a média dos tempos atribuídos a cada execução possa ser computada. Logo, o tempo de execução associado a cada uma das soluções apresentadas no Capítulo \ref{sec:resultados} advém da média dos tempos atribuídos a cinco execuções distintas. 
	
	\item \textbf{Número de avaliações da função objetivo - \boldmath$ n_{ aval} $}
	
	Corresponde ao número de vezes em que a função objetivo foi computada durante a resolução do PCO.
	
	\item \textbf{Máxima violação das restrições - \boldmath$ \Delta c_{max} $ }
	
	Uma vez que o PCO tenha sido resolvido e uma solução tenha sido obtida é necessário confirmar se a solução de fato satisfaz todas as restrições associadas ao PCO. Uma vez que a resolução do PCO consiste em um processo numérico, é inevitável que as restrições sejam levemente violadas. Uma vez computadas as violações associadas a cada restrição, pode-se determinar a violação máxima, utilizada nesse caso para que o atendimento das restrições seja verificado. Considerando que a solução obtida deve satisfazer todas as restrições, espera-se que $ \Delta c_{max} $ seja muito próximo de zero. 
	
	\item \textbf{Número de nós de colocação mínimo - \boldmath$ N_m $}
	
	Com o aumento do número de nós (pontos) de colocação $ N $, normalmente verifica-se a diminuição de $ J^* $. No entanto, a partir de um dado $ N $, denominado número de nós de colocação mínimo, verifica-se que  $ J^* $ se mantém praticamente inalterado à medida que $ N $ cresce. Essa métrica será abordada em detalhes mais adiante.
	
	\item \textbf{Número de execuções bem sucedidas - \boldmath$ N_s $}
	
	Para que a robustez de cada pacote avaliado possa ser verificada, propõe-se que cada estudo de caso seja resolvido considerando-se 30 $ N $ distintos. A cada execução, um $ N $ diferente deve ser adotado, de forma que o número de execuções bem sucedidas indique o quão robusto um dado pacote pode ser. Vale ressaltar que o PPNL advindo do processo de transcrição do PCO depende do número de nós de colocação considerado. No presente trabalho considera-se $ N \in [5, \, N_{max}] $, sendo $ N_{max} $ determinado a partir de um processo de experimentação numérica. 
	
	\item \textbf{Número relativo de execuções bem sucedidas - \boldmath$ N_s\% $}
	
	Consiste na razão entre o número de execuções bem sucedidas e o número total de execuções. Nesse caso $ N_s\% = \dfrac{N_s}{30} $.
\end{enumerate}

\todo[inline, size=normalsize, color=pink]{Falar da métrica Nm}

Observa-se que as métricas 1, 2, 3 e 4 dependem fortemente do número de nós de colocação $ N $ empregado na obtenção de cada solução. Assim sendo, é necessário que um critério para a escolha de $ N $ seja estabelecido. Não seria justo atribuir a todos os métodos o mesmo $ N $, dadas as distintas propriedades numéricas de cada método. Por exemplo, os $ J^* $ associados às soluções obtidas por meio da colocação Hermite-Simpson, costumam ser bem menores que aqueles atribuídos às soluções advindas do emprego da colocação trapezoidal, supondo que o mesmo $ N $ seja utilizado nesses dois casos. De fato, a colocação Hermite-Simpson depende da implementação de nós intermediários, o que melhora a qualidade da solução obtida. Assim sendo, é possível dizer que há métodos que possibilitam a obtenção de melhores resultados a partir de $ N $ menores, capacidade que dificilmente seria avaliada atribuindo-se a todos os métodos o mesmo $ N $. 

Verifica-se que quanto maior o número de nós de colocação empregado na resolução de um PCO, menor tende a ser o $ J^* $ atribuído à solução obtida. No entanto, chega o momento em que o aumento de $ N $ não mais produz reduções significativas em $ J^* $. O número mínimo de nós de colocação é definido de forma que, para $ N > N_m $, não é possível verificar reduções relevantes em $ J^* $. 

Na Figura \ref{fig:metodologia:Nm} é apresentada a relação entre $ N $ e os $ J^* $ atribuídos a dois pacotes quaisquer A e B, sendo $ J^*_{min} $ o valor para o qual converge $ J^* $ quando $ N $ é suficientemente grande. Avaliando-se $ N_{mA} $ e $ N_{mB} $ conclui-se que empregando o pacote B, é possível alcançar $ J^*_{min} $ utilizando um número menor de nós de colocação. Essa característica pode ser considerada uma vantagem do pacote B, dado que normalmente baixos $ N $ estão associados a baixos tempos de processamento \cite{kelly_introduction_2017}. Além disso, conclui-se que utilizar o pacote B adotando-se $ N = N_\infty $ consiste em um desperdício de recursos computacionais.   
 
\noindent	
\begin{minipage}{\textwidth}
	\vspace{\onelineskip}
	\centering
	\includegraphics[width=0.75\linewidth]{draw/metodologia/pdf/Nm}
	\captionof{figure}[Relação entre $ N $ e os $ J^* $ atribuídos a dois pacotes quaisquer A e B]{Relação entre $ N $ e os $ J^* $ atribuídos a dois pacotes quaisquer A e B.}
	\label{fig:metodologia:Nm}
	\vspace{\onelineskip}
\end{minipage}

Para que $ N_m $ seja definido é necessário, primeiramente, que o PCO em análise seja resolvido $ n_m $ vezes, considerando-se, a cada execução, um $ N $ distinto. Em seguida, computa-se $ \overline{J^*} $ a partir da normalização de $ J^* $
%
\begin{equation}
	\overline{J^*} = \frac{J^* - J^*_{min}}{J^*_{max} - J^*_{min}}
\end{equation}
%
em que $ J^*_{min} $ e $ J^*_{max} $ são o menor e o maior valor atribuídos a $ J^* $, de forma que $ 0 \leq \overline{J^*} \leq 1 $. Por fim, defini-se $ N_m $ como sendo o menor $ N $ a partir do qual $ \overline{J} < \epsilon_m $, sendo $ \epsilon_m $ um número próximo de zero. No presente trabalho supõe-se $ n_m = 30 $ e $ \epsilon_m = 0,01 $. O processo de normalização de $ J^* $ e a definição de $ \epsilon_m $, que servem de base para a definição de $ N_m $, são representados na Figura \ref{fig:metodologia:Jnorm}. 

\noindent	
\begin{minipage}{\textwidth}
	\vspace{\onelineskip}
	\centering
	\includegraphics[width=1\linewidth]{draw/metodologia/pdf/Jnorm}
	\captionof{figure}[Representação do processo de normalização de $ J^* $ e da definição de $ \epsilon_m $]{Representação do processo de normalização de $ J^* $ e da definição de $ \epsilon_m $, que servem de base para a definição de $ N_m $.}
	\label{fig:metodologia:Jnorm}
	\vspace{\onelineskip}
\end{minipage}

\todo[inline, size=normalsize, color=pink]{Como estão organizadas as seções do capítulo de resultados}

Os estudos de caso que servem de base para a elaboração do estudo comparativo apresentado no presente trabalho são detalhados nas seções do Capítulo \ref{sec:resultados}. Cada seção se encontra organizada da seguinte maneira:
%
\begin{itemize}
	\item Considerações gerais acerca do estudo de caso em análise;
	\item Formulação da função objetivo, da dinâmica e das restrições associadas ao estudo de caso em análise;
	\item Apresentação dos dados que servem de base para a determinação de $ N_m $;
	\item Análise dos resultados obtidos por meio do emprego de cada um dos pacotes avaliados, considerando-se $ N = N_m $;
	\item Apresentação das trajetórias de estados e controles obtidas a partir da resolução do estudo de caso em análise;
	\item Análise da sensibilidade de $ t_p $ e $ n_{aval} $ ao aumento de $ N $ para cada um dos pacotes avaliados.
\end{itemize} 

Como já foi mencionado anteriormente, o estudo comparativo realizado no presente trabalho é baseado na avaliação de três pacotes: o $ PSOPT $, o $ FALCON $ e o $ COPILOTS $. Uma vez que o $ PSOPT $ e o $ COPILOTS $ possibilitam a utilização de mais de um tipo de colocação, adotam-se os índices $ t $, $ h $ e $ l $ para indicar qual método foi empregado na obtenção de uma dada solução. Esses índices fazem referência à colocação trapezoidal, à colocação Hermite-Simpson, e à colocação pseudo-espectral respectivamente. Para indicar, por exemplo, que uma determinada solução foi obtida empregando-se o $ PSOPT $ e a colocação trapezoidal, atribui-se o rótulo $ PSOPT_t $ à mesma. 

Por fim, ressalta-se que há diversos parâmetros de ajuste associados a cada um dos pacotes avaliados. Assim sendo, seria uma tarefa praticamente impossível verificar a influência de cada um desses parâmetros e ajustá-los de forma que cada execução levasse à melhor solução possível. Logo, recomenda-se que um estudo comparativo se baseei nos valores padrão atribuídos a cada parâmetro \cite{mittelmann_benchmarking_1998}, e, de fato, são esses os valores adotados na obtenção dos resultados apresentados no Capítulo \ref{sec:resultados}.

\todo[inline, color=pink, size=normalsize]{Configurações do PC utilizado}

As configurações do computador utilizado na aquisição e processamento dos dados apresentados no Capítulo \ref{sec:resultados} se encontram listadas a seguir. 
%
\begin{itemize}
	\item \textbf{Modelo}: Laptop Aspire A515-51G;
	\item \textbf{Processador}: Intel\textsuperscript{\textregistered} Core i7-7500U - Frequência máxima de 3,5 GHz - Topologia Dual Core e tecnologia Hyper-Threading Intel\textsuperscript{\textregistered} (2 núcleos e 4 \textit{threads});
	\item \textbf{Memória RAM}: 8 GB (duas placas de 4 GB) DDR4 - Frequência máxima de 2133 MHz; 
	\item \textbf{Placa gráfica}: NVIDIA GeForce 940MX;
	\item \textbf{Sistema operacional}: Linux 64 bits - Kernel 5.8.0-7630 x86\_64;
	\item \textbf{Distribuições}: Pop!\_OS 20.04 LTS ($ FALCON $ e $ COPILOTS $) e Linux Mint 19.3 LTS ($ PSOPT $).
\end{itemize}